% ---------------------------------------------------------------------
%  01-introduction.tex
% ---------------------------------------------------------------------
\section{Introduction}
\label{sec:introduction}

I have spent more than seven years as a senior product manager responsible
for complete products: the hardware, the firmware that runs on it and the
software around it. The portfolio is large and it moves all the time. New
references come in, others are replaced or reach end of life, customers
and sales ask for improvements, competitors force changes, and the product
is expected to get better continuously, not once a year. That kind of job
is less about any single project and more about keeping many product
states, and the work around them, under control at the same time.

Over those years I used every method that product organizations rely on
today: backlogs, Kanban boards, Scrum, roadmaps and Gantt charts,
Stage-Gate reviews, engineering change orders in PLM, and, on the software
side, git and git-flow. Each of them is good at something. None of them
answered the questions I actually had to answer, every week, about my
products.

\subsection{The questions I could not answer}
\label{sec:intro-questions}

Take an illustrative product, the DMS-4000, a fictitious smart door
controller with its own hardware, firmware and companion software. On an
ordinary day its product manager is asked things like these:
\begin{itemize}[leftmargin=1.5em,itemsep=2pt,topsep=3pt]
  \item \emph{A customer reports a defect.} Was it seen on firmware v1.3 or
        on v1.4? Is it already fixed in a later release, or is someone
        still working on it?
  \item \emph{The next release is being planned.} Two improvements and
        three defects are open. Which of them are committed to v1.5? Can
        v1.5 ship if one of them is not finished?
  \item \emph{A component shortage forces a hardware revision.} Which open
        work moves to the new reference? Which stays with the units already
        installed? How long do we keep maintaining the old one?
  \item \emph{An auditor, a key customer or the support team asks} what
        went into v1.4, and why a reported problem was closed without a
        fix.
\end{itemize}

None of these questions is unusual. Yet with the usual tools the answer
means opening a tracker, a spreadsheet, a PLM change order and somebody's
memory, and stitching them together by hand. What is missing is always the
same thing: the link between a piece of work and the \emph{product states}
it belongs to. I kept running into three blind spots:
\begin{enumerate}[label=(\roman*),leftmargin=2.2em,itemsep=1pt,topsep=3pt]
  \item \textbf{Origin.} From which version or hardware reference did this
        work start?
  \item \textbf{Parallelism.} What else is running on the same product, and
        on which line, at the same time?
  \item \textbf{Closure.} Which concrete version incorporated the result,
        or, if none did, why was it abandoned?
\end{enumerate}
Version fields, release milestones and change-order references capture
fragments of this. They are attributes of individual records, not a
structure you can look at, check and plan with. The result is familiar to
anyone in the role: release notes written by hand, defects that cannot be
traced to the version in which they escaped, versions shipped while
committed work was quietly still pending, and hardware revisions whose
relation to the previous reference lives only in people's heads.

\paragraph{When the gap became obvious.}
% AUTHOR: the vignette below is a composite written from the situations
% described above, with no company, product or number. Replace it with a
% real, anonymised story from your own practice if you have one you can
% tell (4-6 sentences, first person, no confidential details).
The gap did not show up in a single project; it showed up as a pattern.
A composite of it, with every detail removed, goes like this. A unit comes
back from the field with a fault, on a hardware reference that had been
revised a few months earlier. Support asks a simple question: is the
firmware fix for this fault already in the build that ships with the new
revision? The fix exists, the ticket is closed and the revision is
released, yet nobody in the room can say whether the one went into the
other without opening three tools and asking two people. Every time I saw
that scene repeat, the missing piece was the same: a record of which
product state each piece of work started from and ended in.

\subsection{What HPB is}
\label{sec:intro-what}

Software engineering solved a structurally identical problem long ago, for
code. In branch-and-merge version control~\citep{chacon2014progit} every
line of work forks from a known commit, runs in parallel with others and
is merged back when it is done; branching models such as
git-flow~\citep{driessen2010gitflow} turned that graph into a picture every
developer recognizes. My proposal is simple: \emph{apply the same model to
the product, not to the code}.

The \emph{Horizontal Product Branch} (HPB) methodology draws each product
as a horizontal main line of baselines (its released and planned
versions). Every improvement, defect, task or laboratory evaluation is a
typed branch that forks from a released baseline and ends in exactly one
of two ways: merged into a named version, or discarded with a written
reason. A hardware revision that creates a new reference forks the line
itself, so the new reference keeps its ancestry. A planned version is a
gate: it does not ship while committed work is still unverified. The name
comes from the drawing: time runs left to right, and the product is the
horizontal line from which all work branches off and to which it returns.

An HPB branch is not a code branch. It is a unit of \emph{change intent}
on the product, which may be implemented by many commits in several
repositories, by a laboratory campaign or by a hardware redesign. HPB
records where that intent started, what happened to it and where it ended.
Throughout the paper, a \emph{product line} means one product reference
followed over time, not a family of variants in the sense of software
product line engineering.

HPB did not come from a single insight. It is the result of accumulating
those blind spots over years of running a real portfolio, trying to patch
them with each tool in turn, and noticing that the patches always pointed
to the same missing structure. I have since built a working prototype of
it inside the product management portal I use, and the improvement points
it captures can be applied well beyond my own products: to any
organization that ships identifiable versions of a product over time,
and especially to products where hardware, firmware and software evolve
together and several references coexist in the field.

\subsection{What this paper gives you}
\label{sec:intro-gives}

This is a methodology paper. It describes a method and everything you need
to apply it consistently; it does not report a controlled study. It
contains:
\begin{enumerate}[label=\textbf{\arabic*.},leftmargin=2em,itemsep=2pt,topsep=3pt]
  \item \textbf{The method.} A small model (product lines, baselines,
        typed branches, merges, discards, line forks and release gates)
        and seven rules, R1--R7, that keep the graph honest
        (\cref{sec:onepage,sec:model}). The full formal model, with
        properties and proof sketches, is in \cref{app:formal}.
  \item \textbf{A standard notation.} Lanes, glyphs and colors with
        redundant encoding, plus layout rules, so that an HPB diagram looks
        the same whoever draws it (\cref{sec:notation}), shown on the
        situations a product manager meets in practice
        (\cref{sec:practice}).
  \item \textbf{An operating guide and metrics.} Practices, cadence, roles,
        coexistence with Scrum, Kanban, Stage-Gate and PLM, an adoption
        path you can start on Monday (\cref{sec:operating}), and flow and
        quality metrics computed directly from the graph
        (\cref{sec:metrics}).
  \item \textbf{Tooling.} HPB Markup, a text language in the spirit of
        Mermaid that turns a few lines of text into an HPB diagram, with an
        open-source renderer for plain HTML, React, Next.js, Vite and the
        command line, and a prototype built into a product management
        portal (\cref{sec:tooling}).
\end{enumerate}
All examples, including the DMS-4000 and ``Product A'', are fictitious and
only illustrate the method. I make no claims about measured effectiveness;
\cref{sec:limitations} is explicit about what has not been shown and how
an organization can evaluate HPB for itself.

\subsection{How to read this paper}
\label{sec:intro-read}

If you have ten minutes, read \cref{sec:gaps} for the problem and
\cref{sec:onepage} for the whole method on one page. If you want to apply
it, add \cref{sec:notation,sec:practice,sec:operating}. If you want to
build tools for it, go to \cref{sec:tooling} and
\cref{app:formal,app:markup,app:package}. \Cref{app:card} is a
quick-reference card of the notation you can print and pin next to your
desk.
